\subsection{Positive offsets and the interior saddle}
\label{sec:direct-positive-offset}

For positive offsets large relative to \(q_1^2\), rescaling reveals
a uniformly hyperbolic interior saddle. Its exponent determines the
single logarithmic term of \(\mathcal L\) near the separatrices, even
as that exponent crosses one. Continuing to the fixed sections adds
a common, possibly unbounded tail multiplier, which cancels from the
endpoint derivative ratio. Away from the saddle, exterior contraction
makes \(\mathcal L\) negative. Together these estimates control all
its zeros on \(\mathcal H_p\).

\begin{proposition}[Positive offset]
\label{prop:direct-positive-offset}
There is a fixed \(\kappa_{\rm p}>0\) such that the return displacement
has at most two distinct zeros when
\[
 q_0>0,\qquad 0\le q_1/\sqrt{q_0}\le\kappa_{\rm p},
\]
for all sufficiently small parameters.
\end{proposition}

\subsubsection{The saddle scale and its local passage}

Put
\[
 \ell=\sqrt{q_0},\qquad \kappa=q_1/\ell,\qquad
 \eta=\ell X,\quad z=\ell^2S,\quad \tau=\ell t .
\]
The scaled field is
\begin{equation}
\label{eq:direct-positive-scaled-field}
 \begin{aligned}
 X_\tau&=1+\kappa X+q_2X^2+\ell X^3R_q(\ell X)
              -SB(\ell X,\ell^2S),\\
 S_\tau&=-XS.
 \end{aligned}
\end{equation}
On every fixed state box it is
\(O_{C^k}(\kappa+|q_2|+\ell)\) close to
\begin{equation}
\label{eq:direct-positive-comparison}
 X_\tau=1-B_0S,\qquad S_\tau=-XS,\qquad B_0=B(0,0,p).
\end{equation}
The comparison field depends on the independent parameters \(A,C\)
through \(B_0\), which lies in a compact positive interval. Its first
integral is
\begin{equation}
\label{eq:direct-positive-first-integral}
 \mathscr H=B_0S-\log S-X^2/2 .
\end{equation}

The equilibrium of the scaled field is \(X=0,\ S=S_s\), where
\(S_sB(0,\ell^2S_s)=1\). It is unique at positive height
in the box in the coordinates \((\eta,z)\), since \(B+zB_z>0\).
Its Jacobian is
\[
 \begin{pmatrix}
 \kappa-\ell S_sB_\eta(0,\ell^2S_s)&
       -(B+\ell^2S_sB_z)(0,\ell^2S_s)\\
 -S_s&0
 \end{pmatrix}.
\]
The positive quantities \(S_s\), \(\lambda_+\) and \(-\lambda_-\), where
\(\lambda_+>0>\lambda_-\) are its eigenvalues, are bounded
above and below uniformly. Moreover,
\begin{equation}
\label{eq:direct-positive-saddle-ratio}
 \nu_p=-\lambda_-/\lambda_+=1+O(\kappa+\ell).
\end{equation}
For passing central heights, the normalized central height is
\(\widehat h_+=h/q_0>S_s\).

Choose small fixed incoming and outgoing sections in the passing
quadrant of this saddle. On each, choose a positive transverse
coordinate equal to \(\mathscr H-\mathscr H_s\) for the comparison
field, and translate its zero to the perturbed separatrix. The
local passage between these coordinates has the form
\begin{equation}
\label{eq:direct-positive-dulac}
 \begin{gathered}
 D_p(x)=u_px^{\nu_p}(1+r_p(x)),\qquad
 u_p=1+O(\kappa+|q_2|+\ell),\\
 |(x\partial_x)^jr_p(x)|\le C_jx^\vartheta,
 \qquad 0\le j\le2,
 \end{gathered}
\end{equation}
for a fixed \(\vartheta>0\). Parameter-uniform expansions of
hyperbolic Dulac maps are developed in \cite{MV2024}. We derive the
finite-order estimate and its leading coefficient below, uniformly as
\(\nu_p\) crosses one.

We first derive the power law on saddle sections where the nonlinear
term is integrable along the passage. We then transfer it to the
chosen sections, keeping the leading coefficient as well as the
Euler remainder.
The parameter-dependent stable manifold theorem
\cite[Theorem~9.6]{Teschl2012}, applied and normalized in
Appendix~\ref{app:positive-saddle-preparation}, gives on a common
saddle box
\begin{equation}
\label{eq:direct-positive-saddle-preparation}
 x'=x,\qquad y'=-\nu_p y+xy^2c_2(x,y).
\end{equation}
The required state derivatives of \(c_2\) are uniformly bounded
for \(\nu_p\) near one.

From input \((x,y_i)\) to output \(x=x_o\), set
\(\tilde y(\xi;x)=y(\xi)(\xi/x)^{\nu_p}\), and write
\[
 k(t,x)=(x/t)^{\nu_p},\qquad
 \mathscr G(t,x,\tilde y)=k(t,x)\tilde y^2c_2(t,k(t,x)\tilde y).
\]
The solution satisfies the Volterra equation
\[
 \tilde y(\xi;x)=y_i+\int_x^\xi
           \mathscr G(t,x,\tilde y(t;x))\,dt .
\]
For \(3/4\le\nu_p\le5/4\) and fixed \(0<\vartheta<1/2\),
\[
 \int_x^{x_o}k(t,x)\,dt
 \le x^{3/4}\int_x^{x_o}t^{-3/4}\,dt
 \le Cx^{3/4}\le Cx^\vartheta .
\]
Fixed small sections therefore keep
\(y_i/2\le\tilde y\le2y_i\). For \(\mathcal E_x=x\partial_x\),
with \(t,\tilde y\) fixed in the integrand, one has
\(\mathcal E_xk=\nu_pk\), hence
\(\lvert\mathcal E_x^i\partial_{\tilde y}^j\mathscr G\rvert
\le C_{i,j}k\) at every required finite order. Euler differentiation
of the Volterra equation gives scalar linear variational equations
with coefficient \(\mathscr G_{\tilde y}\); their forcing terms are
bounded by \(k\) times lower variations. Moving-lower-limit terms
are \(O(x)\): for the first derivative the term is
\(-xy_i^2c_2(x,y_i)\), and another Euler derivative has the same bound.
Variation of constants and the integral estimate give
\[
 |\tilde y(x_o;x)-y_i|
       +\sum_{j=1}^2|\mathcal E_x^j\tilde y(x_o;x)|\le Cx^\vartheta.
\]
Division by the fixed positive \(y_i\) proves the relative Euler
estimates and the leading power on the prepared sections.

The prepared sections and the chosen incoming and outgoing sections
lie in common regular collars at a fixed positive distance from
the saddle. By \cref{lem:regular-transitions}, the transitions
between them, followed by the changes of section coordinates,
contribute positive units, with remainders \(O(x)\) at input and
\(O(x^{\nu_p})\) at output. They preserve the stated Euler estimates. For
\eqref{eq:direct-positive-comparison} the passage in these coordinates
is exactly the identity, by its first integral. Uniform dependence
of the graph constructions, coordinate changes, and fixed section
maps gives \(u_p=1+O(\kappa+|q_2|+\ell)\).
This proves \eqref{eq:direct-positive-dulac}.

The orbits through central points with \(\widehat h_+\downarrow S_s\)
meet the incoming and outgoing sections in their directed order. The
two eigen-coordinates of these central points are positive multiples
of \(\widehat h_+-S_s\), bounded above and below uniformly. Following
backward to the incoming section decreases the unstable coordinate,
so its incoming value tends uniformly to zero. The endpoint order in
\cref{prop:central-displacement-interval-new}, together with the
positive regular section derivative, makes this rebasing injective.
Thus a sufficiently small preliminary interval
\(0<\widehat h_+-S_s<\rho_0\) maps into a fixed collar \(0<x<x_c\).

\subsubsection{Continuation to the fixed sections}

We next continue the maps to the fixed endpoint windows. This step
must include central heights tending to the saddle, whose local passage
time is unbounded. Choose a large fixed \(M_h\), followed by
\(M_t>2M_h+2\), beyond all saddle heights. For
\(S_s<\widehat h_+\le M_h\), both outward graphs reach \(S=M_t\),
with scaled energies \(\mathcal U_\sigma=X_\sigma^2/2
=U_\sigma/q_0\) bounded above and below there. As
\(\widehat h_+\downarrow S_s\), the local construction above makes
the half-orbits converge to the upper separatrices, whose passage from
the fixed saddle box to \(M_t\) is regular.
For the comparison field their energies are exactly
\[
 \mathcal U_\sigma(M_t;\widehat h_+)
   =B_0(M_t-\widehat h_+)-\log(M_t/\widehat h_+),
\]
including the continuous separatrix value at \(\widehat h_+=1/B_0\).
For large \(M_t\) their ratios to \(M_t\) are close to \(B_0\),
uniformly for \(\widehat h_+\le M_h\). Since
\(B_0\in[b_0,\sup B]\) lies in the interior of the common cone, these
ratios lie in that cone.
Outside the saddle box these segments, including the limiting
separatrices, form a compact regular family with uniformly transverse
endpoints at \(S=M_t\). For a small fixed \(\rho_0>0\), use the local
Dulac construction on \(0<\widehat h_+-S_s<\rho_0\), followed by
\cref{lem:regular-transitions} from the saddle sections. On the compact
complement \(S_s+\rho_0\le\widehat h_+\le M_h\), the same lemma applies
from the central section, where the flux has a positive margin.
These continuations preserve the energy bounds and common-cone entry
at \(M_t\). The regular maps have positive bounded first derivatives
and the required higher derivative bounds; the local Dulac map
accounts for the unbounded time.

At the common cut \(t=M_tq_0\),
\[
 q_0/t=1/M_t,\qquad
 q_1/\sqrt t=\kappa/\sqrt{M_t},\qquad
 \sqrt t=\ell\sqrt{M_t}.
\]
Taking \(M_t\) large and the parameters small permits
\cref{lem:direct-weighted-tail} to continue both graphs to \(z_*\).
Keep the tail multiplier in the common endpoint amplitude:
\begin{equation}
\label{eq:direct-positive-amplitude}
 \varpi_{M_tq_0}=(z_*/(M_tq_0))^{-2q_2},\qquad E=q_0\varpi_{M_tq_0} .
\end{equation}
By \eqref{eq:direct-tail-amplitude} with \(t=M_tq_0\),
\(E\le z_*^{1/4}q_0^{3/4}\to0\) as \(q_0\to0\), uniformly for
\(|q_2|<1/8\).

Let \(e_\sigma\) denote the chosen saddle-section coordinates,
with \(e_+=x\) and \(e_-=D_p(x)\). For the comparison field the first
derivatives of the regular maps from these sections to \(S=M_t\) are
both one in the normalized \(I\) coordinate:
at fixed \(S\), a change of \(\mathscr H\) is exactly a change
of \(S-X^2/2\). Their quotient is consequently
\(1+O_{M_t}(\kappa+|q_2|+\ell)\).
The weighted tail gives the endpoint maps
\begin{equation}
\label{eq:direct-positive-endpoints}
 I_\sigma(e)=b_\sigma+E\mathcal F_\sigma(e),\qquad \mathcal F_\sigma(0)=0,\qquad
 c\le\mathcal F_\sigma'\le C,\qquad \|\mathcal F_\sigma\|_{C^3}\le C,
\end{equation}
and
\begin{equation}
\label{eq:direct-positive-slope-quotient}
 \frac{\mathcal F_-'(0)}{\mathcal F_+'(0)}
 \le \exp\left(2C_{\rm cone}\bigl(\sqrt{z_*}+\kappa/\sqrt{M_t}\bigr)\right)
         [1+O_{M_t}(\kappa+|q_2|+\ell)].
\end{equation}
The tail normalization error is
\(O_{C^3}(\ell\sqrt{M_t}+\kappa/\sqrt{M_t})\), and converting from the
cut scale \(M_tq_0\) to \(q_0\) only changes fixed constants. Both tails
start at the common cut, with normalized energies in the common cone
after parameter reduction; their common multiplier \(\varpi_{M_tq_0}\) is
contained in \(E\), and \eqref{eq:direct-tail-first-slope} bounds the
logarithm of each normalized tail slope by
\(C_{\rm cone}(\sqrt{z_*}+\kappa/\sqrt{M_t})\).

If the saddle collar contains a central height in \(\cH_p\), all its
endpoints differ from the admissible endpoints by \(O(E)\), so they lie
in \(\widetilde W_\pm\) for small parameters, by
\cref{lem:exterior-multiplier-new}.

\subsubsection{The return displacement and the remaining central heights}

In the incoming coordinate \(x\), the closing equation is
\[
 \widetilde\Delta_p(x)=\Pi_p(I_-(D_p(x)))-I_+(x)=0.
\]
Its derivative is
\begin{equation}
\label{eq:direct-positive-completed-derivative}
 \widetilde\Delta_p'(x)=E\mathcal F_+'(x)\,[\exp\mathcal L_p(x)-1],
\end{equation}
where
\[
 \mathcal L_p(x)=
 \log \Pi_p'(I_-(D_p(x)))+\log\mathcal F_-'(D_p(x))
             -\log\mathcal F_+'(x)+\log D_p'(x).
\]
The common amplitude \(E\) has disappeared from this ratio. The local
Dulac derivative supplies the logarithmic term, while the endpoint and
exterior factors contribute a bounded constant and a small remainder.
By \eqref{eq:direct-positive-dulac} and
\eqref{eq:direct-positive-slope-quotient},
\begin{equation}
\label{eq:direct-positive-logarithm}
 \mathcal L_p(x)=c_p+(\nu_p-1)\log x+e_p(x),\qquad
 |e_p|+|xe_p'|\le Cx^\vartheta,
\end{equation}
and, after the fixed choices and parameter reductions,
\[
 c_p=\log\left[
 \Pi_p'(b_-)\frac{\mathcal F_-'(0)}{\mathcal F_+'(0)}u_p\nu_p\right]\le-c_0<0.
\]
For the local factor differentiate
\eqref{eq:direct-positive-dulac}: its relative derivative is
\(1+r_p+(x/\nu_p)r_p'\), with the required Euler bound.
The endpoint factors contribute \(O(x)+O(D_p(x))\);
the nonlinear exterior term contributes \(O(ED_p(x))\).
Their first Euler derivatives have the same bounds, by
\eqref{eq:direct-positive-endpoints}. Decrease \(\vartheta\)
if necessary to obtain \eqref{eq:direct-positive-logarithm}.

Using the full constants in \eqref{eq:direct-positive-logarithm}, choose
\(x_c<\min\{1,e^{-1/\vartheta}\}\) such that
\(Cx_c^\vartheta<c_0/2\) and
\(Cx_c^\vartheta|\log x_c|<c_0/4\).
Decrease the preliminary \(\rho_0\) so that its incoming coordinates
lie in \((0,x_c)\); the preceding estimates restrict to this collar.
Use this final \(\rho_0>0\) in the compact complement below, fixing its
continuation constants before the final parameter reduction.
Then \cref{lem:one-logarithm-zero-count} and the negative bound for \(c_p\)
give \(x\mathcal L_p'(x)<0\) at every zero of \(\mathcal L_p\).
To return to the central height, observe that the ratio agrees with
\eqref{eq:direct-completed-log-ratio}: the factor \(dx/dh\) cancels
from the quotient of endpoint derivatives. Moreover,
\(dx/dh=I_{+,h}/(E\mathcal F_+'(x))>0\). Thus the central heights in \(\cH_p\)
with \(0<\widehat h_+-S_s<\rho_0\) form one interval on which
every zero of \(\mathcal L\) has \(\mathcal L_h<0\), that is,
\(\mathcal L_T>0\).

For the remaining central heights we prove contraction. First, for
\(h\ge M_hq_0\),
\[
 |q_0|/h\le1/M_h,\qquad q_1/\sqrt h\le\kappa/\sqrt{M_h}.
\]
For a sufficiently large fixed \(M_h\) and small \(|q_2|\),
\cref{lem:relative-height-paired-new} therefore gives
\(\Delta_h<0\) at these central heights when \(q_1>0\); when
\(q_1=0\), \cref{prop:signed-contraction-new} gives the same sign.
Second, on \(S_s+\rho_0\le\widehat h_+\le M_h\) the radial flux at
the central point has a fixed positive margin. On its compact inner
portion, the scaled linear integral in \eqref{eq:paired-log-derivative-new}
is bounded by \(C_{\rho_0,M_h}\kappa\), and the other nonquadratic integrals
are \(O_{\rho_0,M_h}(\ell)\). To make the tail constant
independent of these height thresholds, continue first to a fixed
large \(S=S_1\). For the comparison field its energy is
\[
 U/q_0=B_0(S-\widehat h_+)-\log(S/\widehat h_+).
\]
Uniformly on \(S_s+\rho_0\le\widehat h_+\le M_h\), its ratio to
\(S\) tends to \(B_0\). Choose \(S_1\) to enter the interior
of the common cone; the positive radial flux at the central point and
finite-height continuation preserve this property for small
\(\kappa,|q_2|,\ell\). The continuation constants depend on the height
thresholds, while the tail constant is fixed by the common cone. By
\eqref{eq:direct-tail-first-slope}, the two tails from \(S=S_1\)
contribute at most \(2C_{\rm cone}(\sqrt{z_*}+\kappa/\sqrt{S_1})\), and
hence
\[
 \left|\log(I_{-,h}/I_{+,h})\right|
 \le 2C_{\rm cone}\sqrt{z_*}
       +C_{\rho_0,M_h,S_1}(\kappa+\ell).
\]
The fixed common-cone choice of \(z_*\), followed by small
\(\kappa_{\rm p},|q_2|,\ell\),
makes this bound smaller than the exterior contraction margin.
These two ranges join into the single interval of large central
heights \(\widehat h_+\ge S_s+\rho_0\), on which \(\mathcal L\) is
strictly negative, including the shared height thresholds.
In increasing \(T\), this interval and the saddle collar, intersected
with \(\cH_p\), cover \(\cH_p\). Since the endpoints of the collar lie in
\(\widetilde W_\pm\), the collar represents this same displacement.
Hence every zero of \(\mathcal L\) on \(\cH_p\) lies in the collar and
crosses upward, and the second case of
\cref{cor:ordered-cover-zero-count} gives at most two zeros of
\(\Delta_p\), proving \cref{prop:direct-positive-offset}.
